%==============================================================================
%  CAPÍTULO IV
%==============================================================================
\chapter{Procedimiento concursal de reorganización judicial ordinario}
\label{cap:reorganizacion-ordinaria}

\fichaCapitulo{Salvataje de la mediana y gran empresa.}{Tramitación de la \pfc{}
ordinaria, quórums y financiamiento preferente.}

%==============================================================================
\bloqueTeorico
%==============================================================================

\subsection{Conservación de la empresa viable frente a la liquidación forzosa}
\pendiente{Desarrollar: efectos de la insolvencia sobre la cadena de proveedores, el
empleo y el sistema financiero.}

\subsection{Naturaleza de la \pfc{} como beneficio extraordinario y temporal}
\pendiente{Desarrollar: alteración de los derechos de cobro individuales del derecho
civil común.}

%==============================================================================
\bloqueNormativo
%==============================================================================

\subsection{Requisitos de admisibilidad de la solicitud}

Análisis del \art{56} \parencite{ley-20720}, con desglose de la declaración jurada de
activos y pasivos.

\subsection{Efectos de la Resolución de Reorganización}

El \art{57} \Nro 1 enumera los efectos de la protección:

\begin{description}[leftmargin=0pt,style=nextline,font=\sffamily\bfseries]
  \item[Letra a) — Inmunidad frente a la ejecución] No puede declararse ni iniciarse un
  procedimiento de liquidación en contra del deudor, ni iniciarse juicios ejecutivos,
  ejecuciones de cualquier clase o restituciones en juicios de arrendamiento.

  \item[Letra b) — Suspensión] Se suspende la tramitación de los procedimientos
  señalados y los plazos de prescripción extintiva.

  \item[Letra c) — Vigencia de los contratos] Todos los contratos suscritos por el deudor
  mantienen su vigencia y condiciones de pago, y no pueden terminarse anticipadamente por
  la sola circunstancia del procedimiento.
\end{description}

\begin{alerta}[La excepción laboral: el límite más importante de la protección]
La inmunidad de la letra a) \destacar{no alcanza a los juicios laborales sobre
obligaciones que gocen de preferencia de primera clase}. Respecto de ellos se suspende
únicamente la ejecución y realización de bienes del deudor, pero el juicio continúa su
tramitación.

La excepción tiene, a su vez, una contraexcepción destinada a impedir el fraude: no opera
cuando se trata de juicios laborales de esa clase que el deudor tuviere a favor de su
cónyuge, de sus parientes, o de los gerentes, administradores, apoderados con poder
general de administración u otras personas con injerencia en la administración de sus
negocios. Para estos efectos se entiende por parientes a los ascendientes, descendientes
y colaterales hasta el cuarto grado de consanguinidad y afinidad, inclusive.

En consecuencia, quien planifica una reorganización debe asumir que el pasivo laboral
preferente seguirá litigándose durante la protección. Véase el capítulo
\ref{cap:creditos-laborales}.
\end{alerta}

\subsection{Duración y prórroga del período de Protección Financiera Concursal}

El \art{57} \Nro 1 fija el plazo base en \destacar{sesenta días} contados desde la
notificación de la Resolución de Reorganización. La \Lreforma{} \citeyearpar{ley-21563}
duplicó el plazo original de treinta días.

El \art{58} regula la prórroga mediante tres vías alternativas, cuyos quórums difieren:

\begin{enumerate}
  \item \textbf{Primera prórroga:} hasta por sesenta días adicionales, con el apoyo de
  dos o más acreedores que representen más del 30\,\% del total del pasivo, excluidos los
  créditos de las personas relacionadas.

  \item \textbf{Segunda prórroga:} hasta por otros sesenta días, solicitada a más tardar
  el décimo día anterior al vencimiento del plazo vigente, con el apoyo de dos o más
  acreedores que representen más del 50\,\% del pasivo, excluidas las personas
  relacionadas.

  \item \textbf{Prórroga única:} en un solo acto, hasta por ciento veinte días, con el
  apoyo de dos o más acreedores que representen más del 50\,\% del pasivo, excluidas las
  personas relacionadas.
\end{enumerate}

El techo es, por cualquiera de los dos caminos, de \destacar{ciento ochenta días}.

\begin{practica}[El plazo de diez días que se pasa por alto]
La segunda prórroga debe solicitarse \emph{hasta el décimo día anterior} al vencimiento
del plazo en curso. No corre desde el vencimiento: corre hacia él. Quien espera a los
últimos días pierde la posibilidad de pedirla.
\end{practica}

\begin{foco}[Los acreedores con garantía real no pierden su preferencia]
El \art{58} precisa que los acreedores hipotecarios y prendarios que apoyen la prórroga
no pierden su preferencia y conservan la facultad de impetrar medidas conservativas. Es
un argumento decisivo al negociar con ellos: apoyar la prórroga no les cuesta su posición.
\end{foco}

\subsection{Financiamiento durante la protección: el crédito preferente}

Régimen de incentivo y preferencia de los créditos contratados durante el período de
protección financiera (\art{74}; \cite{ley-20720}).

%==============================================================================
\bloquePractico
%==============================================================================

\subsection{Iter procesal de la reorganización judicial ordinaria}

\begin{flujograma}{Tramitación de la reorganización judicial ordinaria}{cap04-iter}
  \node[hito] (ingreso) {INGRESO DE SOLICITUD DE REORGANIZACIÓN ORDINARIA};
  \node[nodo, below=of ingreso] (resol)
    {RESOLUCIÓN DE REORGANIZACIÓN (\art{57})\\[2pt]
     \normalfont Inicio de la Protección Financiera Concursal};
  \node[nodo, below=of resol] (verif)
    {VERIFICACIÓN DE CRÉDITOS POR LOS ACREEDORES\\[2pt]
     \normalfont Plazo de 15 días};
  \node[nodo, below=of verif] (informe)
    {INFORME DEL VEEDOR SOBRE LA PROPUESTA DE ACUERDO};
  \node[hito, below=of informe] (junta)
    {CELEBRACIÓN DE LA JUNTA GENERAL DE ACREEDORES};

  \coordinate (split) at ($(junta.south)+(0,-6mm)$);
  \node[favorable, below left=12mm and 4mm of junta.south] (aprueba)
    {APROBACIÓN DEL ACUERDO};
  \node[adverso, below right=12mm and 4mm of junta.south] (rechaza)
    {RECHAZO DE LA PROPUESTA};
  \node[favorable, below=9mm of aprueba] (conv)
    {RESOLUCIÓN QUE TIENE POR APROBADO EL ACUERDO};
  \node[adverso, below=9mm of rechaza] (liq)
    {APERTURA DEL PROCEDIMIENTO DE LIQUIDACIÓN};

  \draw[flecha] (ingreso) -- (resol);
  \draw[flecha] (resol) -- (verif);
  \draw[flecha] (verif) -- (informe);
  \draw[flecha] (informe) -- (junta);
  \draw[flecha] (junta.south) -- (split) -| (aprueba.north);
  \draw[flecha] (junta.south) -- (split) -| (rechaza.north);
  \draw[flecha] (aprueba) -- (conv);
  \draw[flecha] (rechaza) -- (liq);
\end{flujograma}

\subsection{Actuaciones del deudor, del veedor y de la junta}
\pendiente{Desarrollar cada hito con plazos, escritos tipo y errores frecuentes.}

%==============================================================================
\bloqueJurisprudencial
%==============================================================================

\subsection{Límites del control judicial de mérito del acuerdo}
\pendiente{Sistematizar jurisprudencia de la Sala Civil de la \csup{}.}

\subsection{Quitas de capital y acreedores disidentes con garantía real}
\pendiente{Desarrollar el debate doctrinario.}

\subsection{Impugnación por abuso del derecho de acreedores relacionados}
\pendiente{Sistematizar.}
